\documentclass[../../../include/open-logic-section]{subfiles}

\begin{document}

\olfileid{sth}{story}{predicative}	

\olsection{اسنادی تے غیر اسنادی تعریفاں}

رسل دے سیٹ، $R$، دی تعریف $\Setabs{x}{x \notin x}$ نال دتی گئی سی۔
ہور کھول کے دسیے، تاں $R$ اوہ سیٹ ہووے گا جیہدے رکن سارے، تے صرف اوہی،
سیٹ ہوون گے جیہناں بارے ایہ گل سچ نہیں کہ اوہ اپنے آپ دے رکن نہیں۔
سو $R$ دی تعریف وچ اسیں اوس دائرے اُتے مقداری شرط لاندے آں جیہدے وچ
$R$ (جے موجود ہوندا) شامل ہوندا۔

ایہ اک \emph{غیر اسنادی} تعریف اے۔ عام طور اُتے آکھ سکدے آں کہ تعریف
اُدوں تے صرف اُدوں غیر اسنادی اے جدوں اوہ ایسے دائرے اُتے مقداری شرط
لاوے جیہدے وچ تعریف کیتی جاندی شے آپ شامل ہووے۔

پیراڈوکساں دے پچھوں وائٹ ہیڈ، رسل، پوانکارے تے وائل نے ایہو جیہیاں
غیر اسنادی تعریفاں نوں ’’خراب چکر وچ پَون والیاں‘‘ آکھ کے ردّ کیتا:
\begin{quote}
	جیہڑے پیراڈوکس توں بچنا اے، اوہناں دی پڑتال توں پتا لگدا اے کہ سارے
	اک طرح دے خراب چکر توں نکلدے نیں۔ ایہ خراب چکر ایس مفروضے توں پیدا
	ہوندے نیں کہ چیزاں دے کسے اکٹھ وچ کجھ ایہو جیہے رکن ہو سکدے نیں
	جیہناں دی تعریف صرف پورے اکٹھ دے ذریعے ہو سکدی اے[\ldots. \textparagraph]

	جس اصول نال اسیں ناجائز کُلّاں توں بچ سکدے آں، اوہ ایویں بیان ہو
	سکدا اے۔ ’کسے اکٹھ دے \emph{سارے} رکن ورتن والی کوئی شے اوس اکٹھ دی
	رکن نہیں ہو سکدی‘۔ یا اُلٹے پاسے، ’جے کوئی اکٹھ کُلّ ہون دی صورت وچ
	ایہو جیہے رکن رکھدا ہووے جیہڑے صرف اوس کُلّ نال تعریف ہو سکن، تاں
	اوس اکٹھ دا کوئی کُلّ نہیں۔‘ اسیں ایہنوں ’خراب چکر دا اصول‘ آکھاں گے،
	کیوں جے ایہ سانوں ناجائز کُلّ منّن نال پیدا ہون والے خراب چکراں توں
	بچاندا اے۔
	\citep[p.~37]{WhiteheadRussell1910}
\end{quote}
جے اسیں \emph{خراب چکر دا اصول} ردّ کرن وچ اوہناں دے پچھے چلیے، تاں اسیں
(\olref[sth][story][rus]{sec} والے) تباہ کن سادے خاصیت توں سیٹ بناؤن دے
خاکے دی تھاں ایہو جیہا اصول رکھن دی کوشش کر سکدے آں:

\begin{defish}
\emph{اسنادی طور اُتے خاصیت توں سیٹ بنانا۔} ہر اوہ فارمولا $\phi$ جیہڑا صرف سیٹاں اُتے مقداری شرط لاندا اے، اوہدے لئی سیٹ$^\prime$ $\Setabs{x}{\phi(x)}$ موجود اے۔
\end{defish}

جِنّی دیر سیٹ$^{\prime}$ سیٹ نہیں، کوئی تضاد نہیں نکلے گا۔

بدقسمتی نال اسنادی طور اُتے خاصیت توں سیٹ بناؤن دا اصول بہتا \emph{وسیع}
نہیں۔ ایہ سانوں نویں قسم دیاں شےواں، سیٹ$^\prime$، نال ملاندا اے۔
سو سانوں اوہ فارمولے وی ویکھنے پین گے جیہڑے سیٹ$^\prime$ اُتے مقداری شرط
لاندے نیں۔ جے اوہ ہمیشہ سیٹ$^\prime$ دین، تاں رسل دا پیراڈوکس دوبارہ
پیدا ہو جائے گا۔ صرف اوہ سیٹ$^\prime$ ویکھ لو جیہدے وچ سارے اوہ سیٹ$^\prime$ ہون
جیہڑے اپنے آپ دے رکن نہیں۔ سو ایس خیال نوں اگے لے جائیے تاں آکھنا پئے گا
کہ سیٹ$^\prime$ اُتے مقداری شرط لاؤن والا فارمولا اوہدے مطابق اک سیٹ$^{\prime\prime}$
دیندا اے۔ فیر سانوں سیٹ$^{\prime\prime\prime}$، سیٹ$^{\prime\prime\prime\prime}$،
وغیرہ چاہیدے ہون گے۔ اینے سارے پرائم نشان توں بچن لئی ایہناں نوں سیٹ$_0$،
سیٹ$_1$، سیٹ$_2$، سیٹ$_3$، سیٹ$_4$، \ldots سمجھنا سوکھا اے۔
ایس طرح (سادے) قسماں دے نظریے تک اک راہ مل جاندی اے۔

ایہو جیہے نظریے اُتے کجھ ظاہر اعتراض نیں (پر ایہ ظاہر نہیں کہ ایہ
\emph{فیصلہ کن} اعتراض نیں)۔ مختصر گل ایہ اے کہ نتیجے والا نظریہ ورتنا
اوکھا اے، ایہ کئی وکھ وکھ قسماں دیاں شےواں فراخ دلی نال منّ لیندا اے،
تے آخر ایہ وی صاف نہیں کہ غیر اسنادی تعریفاں \emph{ایڈیاں ماڑیاں} نیں۔

آخری گل صاف کرن لئی \citeauthor{Ramsey1925} دا ایہ قول ویکھو:
\begin{quote}
	اسیں کسے بندے نوں اک ٹولی دا سبھ توں لمّا بندہ آکھ سکدے آں۔ ایویں
	اوہ اک ایسے کُلّ دے ذریعے پچھانیا جاندا اے جیہدا اوہ آپ رکن اے،
	تے کوئی خراب چکر وی نہیں پیندا۔ \citep{Ramsey1925}
\end{quote}
ریمزی دی گل ایہ اے کہ ’’ٹولی دا سبھ توں لمّا بندہ‘‘ \emph{واقعی} اک غیر اسنادی
تعریف، پر ایہ صاف طور اُتے بالکل جائز اے۔

کوئی ایہ جواب دے سکدا اے کہ ایس صورت وچ اسیں سبھ توں لمّے بندے دی پچھان
\emph{اسنادی} طریقے نال وی کرا سکدے آں۔ مثلاً اسیں اوس بندے ول صرف اشارہ
کر سکدے آں۔ فیر غیر اسنادی تعریفاں اُتے اعتراض صرف اوہناں شےواں تک
رکھنا پئے گا جیہناں دی پچھان \emph{صرف} غیر اسنادی طریقے نال ہو سکدی اے۔
پر فیر وی ایہ ہور سمجھنا پئے گا کہ ایہو جیہی ’’لازمی غیر اسنادیت‘‘
کاہتوں اینی ماڑی اے۔\footnote{ہور تفصیل لئی \citet{Linnebo2010} ویکھو۔}

ایہ منّن والی گل اے کہ جے اسیں تعریف توں ایہ چاہندے آں کہ اوہ سانوں
کوئی شے \emph{بنان} دا طریقہ دسے، تاں غیر اسنادی تعریفاں بڑی اوکڑ پاندیاں
نیں۔ ایہو جیہی تعریف دتی ہووے تاں بندہ سچ مچ خراب چکر وچ پھس جاندا اے۔
غیر اسنادی طور اُتے دسی شے بناؤن لئی \emph{پہلاں} ساریاں شےواں بناؤنیاں
پین گیاں (اوہ غیر اسنادی طور اُتے دسی شے وی)، کیوں جے اوہ شے ساریاں شےواں
دے لحاظ نال دسی گئی اے۔ سو اوہ شے بناؤن توں پہلاں اوہنوں آپ ای بنانا پئے گا۔
پر ’’لازمی طور اُتے غیر اسنادی‘‘ طریقے نال دسے سیٹاں خلاف ایہ سنجیدہ اعتراض
صرف اوہدوں اے جے اسیں سیٹاں نوں ایہو جیہیاں شےواں سمجھدے آں جیہڑیاں
اسیں \emph{بناندے} آں۔ تے اسیں (شاید) ایویں نہیں سمجھدے۔

ایس طرح، چنگا ہویا یا ماڑا، عام ہویا طریقہ (غیر)اسنادیت بارے سخت پینتڑا
نہیں لیندا۔ ایہدے وچ اوہ طریقہ ورتیا جاندا اے جیہڑا ہُن اکٹھا کرن والا
مرحلہ وار طریقہ سمجھیا جاندا اے۔ آخر ایہ سانوں اپنے سیٹ ’’مرحلیاں‘‘ وچ
ونڈن دیوے گا۔ ایہ \emph{کجھ کجھ} اویویں اے جیویں اسنادی طریقہ شےواں نوں
سیٹ$_0$، سیٹ$_1$، سیٹ$_2$، \ldots وچ ونڈدا اے۔ پر اسیں اوہناں دی قسم وچ
کوئی فرق فرض نہیں کراں گے۔

\end{document}
